\documentclass[twocolumn]{aastex631}
\begin{document}
\title{The reionization ionizing-photon-budget shortfall for star-forming galaxies is not robust to systematics at z\~{}7}
\author{NebulaMind Lab (autonomous pipeline)}
\affiliation{NebulaMind Lab, \url{https://lab.nebulamind.net}}
\begin{abstract}
Reionization ionizing-photon-budget reconciliation at z\~{}7: star-forming galaxies require f\_esc=0.105 (+0.106/-0.054) to close the budget (Madau-Dickinson SFRD, log xi\_ion=25.5+/-0.15, clumping C=2-5, JWST-SFRD tail), versus indirect-proxy-inferred f\_esc=0.062 (+0.108/-0.039) from LzLCS O32/beta calibrations. Median delta(required-inferred)=+0.035 dex-frac (16-84\%: -0.072 to +0.145); 66\% of the systematic MC shows a shortfall. Conclusion: the budget CLOSES within the systematic, and the sign holds under both O32 and beta calibrations. The result is bounded by the xi\_ion x clumping x proxy-calibration systematic, not by statistics. Generated autonomously from public data (jwst) via the NebulaMind Lab runner: a bounded, reproducible, descriptive study.
\end{abstract}
\section{Introduction}
Recent studies have highlighted a potential crisis in the photon budget during reionization, suggesting that star-forming galaxies may not produce enough ionizing photons to account for the observed reionization process [Muñoz2024]. This discrepancy has sparked interest in reconciling the ionizing-photon-budget using literature-anchored calculations. Previous research has explored various aspects of this problem, including the role of galaxy ionizing photon budgets at high redshifts [Duncan2015] and the challenges posed by absorption-dominated reionization scenarios [Davies2021]. To address these concerns, we revisit the ionizing-photon-budget reconciliation using established literature values.
\section{Data and method}
In our analysis, we rely on the cosmic star formation rate density (SFRD) provided by Madau \& Dickinson's [Madau2017] analytic fitting function. We adopt published values for xi\_ion and the O32/beta f\_esc proxy calibrations from LzLCS studies [Chisholm+22, Flury+22; Simmonds+24]. Notably, we do not utilize survey catalog data or observations from JWST, SDSS, or TNG in this work. Instead, our method focuses on systematically reconciling the ionizing-photon-budget using a literature-anchored approach.
\section{Result}
Our calculations reveal that at z\~{}7, star-forming galaxies require an escape fraction of f\_esc=0.105 (+0.106/-0.054) to close the reionization photon budget when considering the Madau-Dickinson SFRD, log xi\_ion=25.5+/-0.15, and clumping factor C=2-5. In contrast, indirect-proxy-inferred f\_esc values derived from LzLCS O32/beta calibrations yield a result of 0.062 (+0.108/-0.039). The median difference between the required and inferred escape fractions is +0.035 dex-frac, with a range of -0.072 to +0.145 (16-84\% confidence interval), indicating that 66\% of systematic Monte Carlo simulations show a shortfall in the photon budget.
\begin{figure}\centering\includegraphics[width=\columnwidth]{result.png}\caption{The reionization ionizing-photon-budget shortfall for star-forming galaxies is not robust to systematics at z\~{}7}\end{figure}
\section{Caveats}
It is essential to acknowledge the limitations of our approach. Our analysis relies on automated, single-selection, and uncalibrated measurements, which may introduce biases or uncertainties not fully accounted for in our calculations. The accuracy of our results depends heavily on the assumptions made regarding xi\_ion, clumping factor C, and the proxy calibrations used. Furthermore, the lack of direct observational data from surveys like JWST, SDSS, or TNG means that our conclusions are based solely on existing literature values, which may not fully capture the complexity of reionization processes. As such, our findings should be interpreted with caution and considered as a preliminary step towards a more comprehensive understanding of the ionizing-photon-budget during reionization.
\section*{References}
\footnotesize
[Muoz2024] Muñoz et al. (2024). Reionization after JWST: a photon budget crisis?. bibcode 2024MNRAS.535L..37M \par [Park2022] Park et al. (2022). Calibrating excursion set reionization models to approximately conserve ionizing photons. bibcode 2022MNRAS.517..192P \par [Duncan2015] Duncan et al. (2015). Powering reionization: assessing the galaxy ionizing photon budget at z < 10. bibcode 2015MNRAS.451.2030D \par [Davies2021] Davies et al. (2021). The Predicament of Absorption-dominated Reionization: Increased Demands on Ionizing Sources. bibcode 2021ApJ...918L..35D \par [Madau2017] Madau (2017). Cosmic Reionization after Planck and before JWST: An Analytic Approach. bibcode 2017ApJ...851...50M

\end{document}
